

% \noindent \red{\textbf{ASSIGNMENTS:}}

% \textbf{XIANGGUO} in charge of:
% \begin{itemize}
%     \item 
% \end{itemize}

% \jiawen{\textbf{JIAWEN in charge of:}
% \begin{itemize}
%     \item 
% \end{itemize}
% }

% \hrule\vspace*{\baselineskip} 




\noindent \red{\textbf{TIMELINE: (Aug.-Oct. Clear within 3 Months)}}
{\color{gray}{
\begin{enumerate}
    \item Aug: collect and analyze papers; start to prepare for section 7.4 (Xiangguo will contact Mr. Chenyi ZI for further details)
    \item Sep: 1st draft finished, open-source finished
    \item Oct: 2nd version finished, open-source polished, submit by the end of Oct. 
\end{enumerate}
}}

\hrule\vspace*{\baselineskip} 

\noindent \red{\textbf{ROADMAP:}}

\noindent \%cutline of the colors:
\begin{itemize}
    \item \jwxg{work done by}
    \item \jiawen{work done by}
    \item work done by (Xiangguo)
    \item \red{Red Color: highlight of our survey}
\end{itemize}


\begin{enumerate}[1]
    \item Introduction
    \item Survey Methodology
    \begin{enumerate}[2.1]
        \item Research Questions
        \item Taxonomy
        \item Literature Overview
    \end{enumerate}
\item Preliminaries
\begin{enumerate}[3.1]
    \item \jiawen{Graph Neural Networks}
    \item \jiawen{Pre-training, and Fine-tuning}
    \item \jiawen{A Brief History of Prompt Learning}
    \item \red{Why Prompt? A New Perspective upon Flexibility and Expressiveness.}: \note{Xiangguo in charge of this}
\end{enumerate}
\item \jiawen{Pre-training Methods on Graph Domains}
\begin{enumerate}[4.1]
    \item \jiawen{node-level strategies}
    \item \jiawen{edge-level strategies}
    \item \jiawen{graph-level strategies}
    \item \jiawen{multi-task pre-training}
    \item \jwxg{Further Discussion:} \\\note{Jiawen writes the content. Xiangguo supplements anything insightful here. } limitation and discussion, correlation to the prompting performance. 
\end{enumerate}
\item \jwxg{Prompting Methods for Graph Models}\\
\note{Jiawen help to find ``food'', and supplement anything insightful; Xiangguo writes the main content}
\begin{enumerate}[5.1]
    \item Prompt Design for Graph Tasks
    \begin{enumerate}[5.1.1]
        \item [5.1.1] Graph Prompt for Node-level Tasks
        \item [5.1.2] Graph Prompt for Edge-level Tasks
        \item [5.1.3] Graph Prompt for Graph-level Tasks
        \item [5.1.4] Multi-task Prompting on Graphs
    \end{enumerate}
    \item Multi-model Prompting using Graphs: \\\note{TBD: this section can be also moved to the Application section}
    \begin{enumerate}[5.2.1]
        \item Text-Graph: \note{also include knowledge graph}
        \item More: \note{multi-modal prompting with graph}
    \end{enumerate}
    \item \red{Discussion in the Unified Framework}: 
    \note{I will analyze the above work using our proposed framework: prompt tokens, token structure, inserting patterns} 
\end{enumerate}
\item \jiawen{Applications}
\\\note{No need to find prompt-based work (it would be better if you can find any prompt-based work, although), just discuss potential applications that might be integrated with a prompt in the future} e.g.: Medicine (TBD),  Biology (TBD), Knowledge management (TBD), Online Social Networks (TBD)... Each application corresponds to one subsection
\item \jwxg{Discussion}
\\\note{Jiawen supplements anything insightful here; Jiawen finishes section 7.1 and 7.4; Xiangguo writes the rest and polishes the whole content}
\begin{enumerate}[7.1]
    \item \jiawen{Summary of this paper}
    \item Current Challenges: each paragraph to each challenge
    \item Future Directions
    \begin{enumerate}[7.3.1]
        \item finding knowledge from large graph models (LGMs) like LLMs: we are still waiting for the era of LGMs because currently, the parameters of most graph models are far from heavy than LLMs.
        \item using a prompt for more transferable learning.
        \item well theoretical work on the effectiveness of graph prompt
        \item multi-task pre-training to multi-task prompting
        \item explainable and understandable prompt design for graph
    \end{enumerate}
    \item \jiawen{Resources }
    \begin{enumerate}[7.4.1]
        \item \jiawen{ProG: A Unified Library for Graph Prompting}
        \item \jiawen{Website collecting literature, datasets, and available codes.}
    \end{enumerate}
\end{enumerate}
\item Conclusion
\end{enumerate}
\clearpage

